\section{Conclusions}\label{sec:conclusions}

We have described a quintic B-spline collocation method for the continuously monitored arithmetic Asian
call under the Heston model and, with most of the effort, the references needed to check it. The method
is simple to state. It solves the two-dimensional equation obtained from the Rogers--Shi--Ve\v{c}e\v{r}
reduction; the time-independence of that equation lets the whole time march run on two sparse
factorisations; the degenerate boundary $v=0$ is handled by collocating the equation on it; and the price
at any strike is an exact evaluation of the spline.

Against a Monte Carlo reference with a standard error of about $1.6\times10^{-5}$, built from the lower
bound of \citet{FusaiKyriakou2016} used as a control variate, the solver agrees within the reference's
uncertainty at eleven strikes. Over thirteen parameter regimes, including a Feller ratio of $0.08$, the
differences are as large as that uncertainty leads one to expect ($\chi^2=45.6$ for 39 comparisons), once
the domain has been widened where the regime requires it. The
time discretisation is of second order, with a cubic remainder from the Rannacher start-up that two levels
of Richardson extrapolation remove. In space the observed order approaches four from below, the limit of
collocation at Greville points, and the error at a moderate grid of 273 points in $x$ is below
$10^{-6}$. On the same equation and at equal computing time, a modified Craig--Sneyd ADI scheme has an error
about 150 times larger, and within 25 seconds it does not get below $5\times10^{-3}$.

Some results are useful beyond this particular solver. The lower bound of Fusai and Kyriakou can be
computed for the Heston model with a continuous average by two consecutive Riccati integrations, and once
computed it makes an unusually effective control variate. Truncation errors of the computational domain
can be of order $10^{-4}$ while every refinement study looks clean; they showed up only against an
independent reference, at vol-of-vol $\xi=1.5$ and at a five-year maturity. And the published lower bounds
for the Heston cases in \citet{FusaiKyriakou2016} are not consistent with our computations of the same
quantity for the model as printed; we have listed what we could and could not rule out.

The work has limits. We considered one model, one contract and no dividends. The references all rely on
the affine structure of Heston; for a non-affine model the checking strategy would have to change. Our
spatial order is capped at four by the choice of collocation points. For maximally smooth splines there
are two known ways around this cap, moving the points to Galerkin-superconvergent positions
\citep{Anitescu2015,Montardini2017} or keeping them and adding a high-order perturbation of the operator
\citep{IrodotouEllinaHoustis1988,ChristaraNg2006}; classical Gauss-point collocation
\citep{deBoorSwartz1973,BialeckiFairweather2001} would require a less smooth spline space and would lose
the property that the degenerate boundary is a collocation line. None of these has been analysed for a
degenerate operator with a mixed derivative and a kinked terminal condition, and at the accuracy that our
references can confirm there is little to gain yet. Other natural extensions are discrete monitoring,
sensitivities, which the spline representation gives directly, and floating-strike contracts, which the
same reduction covers.
